% Capitulo 4: matematicas con referencias cruzadas y derivacion.
\section{Matemáticas}\label{sec:matematicas}

La composición matemática es la de \LaTeX{} de siempre, con
\texttt{amsmath} cargado en el preámbulo. En línea, la fórmula
cuadrática $x = \frac{-b \pm \sqrt{b^2 - 4ac}}{2a}$ aparece dentro del
párrafo sin romper el ritmo, y la identidad de Euler $e^{i\pi} + 1 = 0$
recuerda que cinco constantes caben en una línea. La varianza de una
muestra, $\sigma^2 = \frac{1}{n}\sum_{i=1}^{n}(x_i - \mu)^2$, conecta este
capítulo con el guion de estadística del capítulo~\ref{sec:codigo}.

\subsection{Ecuaciones numeradas}

La integral de Gauss tiene valor exacto conocido, recogido en la
ecuación~\eqref{eq:gauss}. Su demostración clásica eleva la integral al
cuadrado y cambia a coordenadas polares; el resultado, $\sqrt{\pi}$,
reaparece cada vez que una normal necesita normalizarse.

\begin{equation}\label{eq:gauss}
  \int_{-\infty}^{\infty} e^{-x^2}\,dx = \sqrt{\pi}
\end{equation}

Junto a ella, la fórmula de Euler relaciona la exponencial compleja con el
seno y el coseno, como muestra la ecuación~\eqref{eq:euler}. Entre la
ecuación~\eqref{eq:gauss} y la ecuación~\eqref{eq:euler} queda
demostrado que las referencias cruzadas con \texttt{\textbackslash eqref}
funcionan: cada número entre paréntesis enlaza con su ecuación y
\texttt{texforge check} protesta si la etiqueta no existe.

\begin{equation}\label{eq:euler}
  e^{ix} = \cos x + i\,\sin x
\end{equation}

\subsection{Una derivación}

La media aritmética minimiza la suma de desvíos al cuadrado, y la
derivación cabe en un entorno \texttt{align} con una sola línea
numerada. Derivar respecto de $\mu$ e igualar a cero da directamente el
estimador, como se ve en la ecuación~\eqref{eq:derivacion}.

\begin{align}
  S(\mu) &= \sum_{i=1}^{n}(x_i - \mu)^2 \nonumber \\
  \frac{dS}{d\mu} &= -2\sum_{i=1}^{n}(x_i - \mu) \nonumber \\
  \hat{\mu} &= \frac{1}{n}\sum_{i=1}^{n}x_i \label{eq:derivacion}
\end{align}

La ecuación~\eqref{eq:derivacion} es la misma media que calcula la
función \texttt{mean} del listado~\ref{lst:archivo}: la teoría de este
capítulo y el código del capítulo~\ref{sec:codigo} cuentan la misma
historia por dos caminos, y ambos llegaron al PDF desde texto plano.

\subsection{Matrices}

Los vectores y las matrices se componen con los entornos de amsmath, y las
referencias funcionan igual que con las ecuaciones: la matriz identidad de
la ecuación~\eqref{eq:identidad} reaparece cada vez que un sistema lineal
necesita su elemento neutro.

\begin{equation}\label{eq:identidad}
  \mathbf{I} = \begin{pmatrix}
    1 & 0 & \cdots & 0 \\
    0 & 1 & \cdots & 0 \\
    \vdots & \vdots & \ddots & \vdots \\
    0 & 0 & \cdots & 1
  \end{pmatrix}
\end{equation}

La ecuación~\eqref{eq:identidad} multiplica sin cambiar, como el uno de la
aritmética: para toda matriz cuadrada vale que el producto con la identidad
devuelve la matriz intacta. Entre la integral de Gauss de la
ecuación~\eqref{eq:gauss} y esta identidad queda cubierto lo esencial:
integrales, exponenciales complejas, derivaciones alineadas y álgebra lineal
conviven en el mismo documento sin paquetes adicionales.
